\section{Trace and norm}
\label{section-trace-pairing}

\noindent
Let $L/K$ be a finite extension of fields. By
Lemma \ref{lemma-vector-space-is-free}
we can choose an isomorphism $L \cong K^{\oplus n}$ of $K$-modules.
Of course $n = [L : K]$ is the degree of the field extension.
Using this isomorphism we get for a $K$-algebra map
$$
L \longrightarrow \text{Mat}(n \times n, K),\quad
\alpha \longmapsto \text{matrix of multiplication by }\alpha
$$
Thus given $\alpha \in L$ we can take the trace and the determinant
of the corresponding matrix. Of course these quantities are independent
of the choice of the basis chosen above. More canonically, simply thinking
of $L$ as a finite dimensional $K$-vector space we have
$\text{Trace}_K(\alpha : L \to L)$ and the determinant
$\det_K(\alpha : L \to L)$.

\begin{definition}
\label{definition-trace-norm}
Let $L/K$ be a finite extension of fields. For $\alpha \in L$ we define
the {\it trace}
$\text{Trace}_{L/K}(\alpha) = \text{Trace}_K(\alpha : L \to L)$
and the {\it norm}
$\text{Norm}_{L/K}(\alpha) = \det_K(\alpha : L \to L)$.
\end{definition}

\noindent
It is clear from the definition that
$\text{Trace}_{L/K}$ is $K$-linear and satisfies
$\text{Trace}_{L/K}(\alpha) = [L : K]\alpha$ for $\alpha \in K$.
Similarly $\text{Norm}_{L/K}$ is multiplicative and
$\text{Norm}_{L/K}(\alpha) = \alpha^{[L : K]}$ for $\alpha \in K$.
This is a special case of the more general construction discussed
in Exercises, Exercises \ref{exercises-exercise-trace-det} and
\ref{exercises-exercise-trace-det-rings}.

\begin{lemma}
\label{lemma-characteristic-vs-minimal-polynomial}
Let $L/K$ be a finite extension of fields. Let $\alpha \in L$ and let $P$
be the minimal polynomial of $\alpha$ over $K$. Then the characteristic
polynomial of the $K$-linear map $\alpha : L \to L$ is equal to
$P^e$ with $e \deg(P) = [L : K]$.
\end{lemma}

\begin{proof}
Choose a basis $\beta_1, \ldots, \beta_e$ of $L$ over $K(\alpha)$.
Then $e$ satisfies $e \deg(P) = [L : K]$ by
Lemmas \ref{lemma-degree-minimal-polynomial} and
\ref{lemma-multiplicativity-degrees}.
Then we see that $L = \bigoplus K(\alpha) \beta_i$ is a
direct sum decomposition into $\alpha$-invariant subspaces
hence the characteristic polynomial of $\alpha : L \to L$
is equal to the characteristic polynomial of
$\alpha : K(\alpha) \to K(\alpha)$ to the power $e$.

\medskip\noindent
To finish the proof we may assume that $L = K(\alpha)$.
In this case by Cayley-Hamilton we see that $\alpha$
is a root of the characteristic polynomial. And since the
characteristic polynomial has the same degree as the minimal
polynomial, we find that equality holds.
\end{proof}

\begin{lemma}
\label{lemma-trace-and-norm-from-minimal-polynomial}
Let $L/K$ be a finite extension of fields. Let $\alpha \in L$ and let
$P = x^d + a_1 x^{d - 1} + \ldots + a_d$
be the minimal polynomial of $\alpha$ over $K$. Then
$$
\text{Norm}_{L/K}(\alpha) = (-1)^{[L : K]} a_d^e
\quad\text{and}\quad
\text{Trace}_{L/K}(\alpha) = - e a_1
$$
where $e d = [L : K]$.
\end{lemma}

\begin{proof}
Follows immediately from Lemma \ref{lemma-characteristic-vs-minimal-polynomial}
and the definitions.
\end{proof}

\begin{lemma}
\label{lemma-trace-and-norm-linear}
Let $L/K$ be a finite extension of fields. Let $V$ be a finite dimensional
vector space over $L$. Let $\varphi : V \to V$ be an $L$-linear map.
Then
$$
\text{Trace}_K(\varphi : V \to V) =
\text{Trace}_{L/K}(\text{Trace}_L(\varphi : V \to V))
$$
and
$$
\det\nolimits_K(\varphi : V \to V) =
\text{Norm}_{L/K}(\det\nolimits_L(\varphi : V \to V))
$$
\end{lemma}

\begin{proof}
Choose an isomorphism $V = L^{\oplus n}$ so that $\varphi$ corresponds
to an $n \times n$ matrix. In the case of traces, both sides of the formula
are additive in $\varphi$. Hence we can assume that $\varphi$
corresponds to the matrix with exactly one nonzero entry in the $(i, j)$ spot.
In this case a direct computation shows both sides are equal.

\medskip\noindent
In the case of norms both sides are zero if $\varphi$ has a nonzero kernel.
Hence we may assume $\varphi$ corresponds to an element of
$\text{GL}_n(L)$. Both sides of the formula are multiplicative in $\varphi$.
Since every element of $\text{GL}_n(L)$ is a product of elementary
matrices we may assume that $\varphi$ either looks like
$$
E_{12}(\lambda) =
\left(
\begin{matrix}
1 & \lambda & \ldots \\
0 & 1 & \ldots \\
\ldots & \ldots & \ldots
\end{matrix}
\right)
\quad\text{or}\quad
E_1(a) =
\left(
\begin{matrix}
a & 0 & \ldots \\
0 & 1 & \ldots \\
\ldots & \ldots & \ldots
\end{matrix}
\right)
$$
(because we may also permute the basis elements if we like).
In both cases the formula is easy to verify by direct computation.
\end{proof}

\begin{lemma}
\label{lemma-trace-and-norm-tower}
Let $M/L/K$ be a tower of finite extensions of fields. Then
$$
\text{Trace}_{M/K} = \text{Trace}_{L/K} \circ \text{Trace}_{M/L}
\quad\text{and}\quad
\text{Norm}_{M/K} = \text{Norm}_{L/K} \circ \text{Norm}_{M/L}
$$
\end{lemma}

\begin{proof}
Think of $M$ as a vector space over $L$ and apply
Lemma \ref{lemma-trace-and-norm-linear}.
\end{proof}

\noindent
The trace pairing is defined using the trace.

\begin{definition}
\label{definition-trace-pairing}
Let $L/K$ be a finite extension of fields. The {\it trace pairing}
for $L/K$ is the symmetric $K$-bilinear form
$$
Q_{L/K} : L \times L \longrightarrow K,\quad
(\alpha, \beta) \longmapsto \text{Trace}_{L/K}(\alpha\beta)
$$
\end{definition}

\noindent
It turns out that a finite extension of fields is separable if and only
if the trace pairing is nondegenerate.

\begin{lemma}
\label{lemma-separable-trace-pairing}
Let $L/K$ be a finite extension of fields. The following are equivalent:
\begin{enumerate}
\item $L/K$ is separable,
\item $\text{Trace}_{L/K}$ is not identically zero, and
\item the trace pairing $Q_{L/K}$ is nondegenerate.
\end{enumerate}
\end{lemma}

\begin{proof}
It is clear that (3) implies (2). If (2) holds, then pick $\gamma \in L$
with $\text{Trace}_{L/K}(\gamma) \not = 0$. Then if $\alpha \in L$
is nonzero, we see that $Q_{L/K}(\alpha, \gamma/\alpha) \not = 0$.
Hence $Q_{L/K}$ is nondegenerate. This proves the equivalence of
(2) and (3).

\medskip\noindent
Suppose that $K$ has characteristic $p$ and $L = K(\alpha)$ with
$\alpha \not \in K$ and $\alpha^p \in K$. Then $\text{Trace}_{L/K}(1) = p = 0$.
For $i = 1, \ldots, p - 1$ we see that $x^p - \alpha^{pi}$ is the minimal
polynomial for $\alpha^i$ over $K$ and we find
$\text{Trace}_{L/K}(\alpha^i) = 0$ by
Lemma \ref{lemma-trace-and-norm-from-minimal-polynomial}.
Hence for this kind of purely inseparable degree $p$ extension
we see that $\text{Trace}_{L/K}$ is identically zero.

\medskip\noindent
Assume that $L/K$ is not separable. Then there exists a subfield
$L/K'/K$ such that $L/K'$ is a purely inseparable degree $p$ extension
as in the previous paragraph, see
Lemmas \ref{lemma-separable-first} and \ref{lemma-finite-purely-inseparable}.
Hence by Lemma \ref{lemma-trace-and-norm-tower}
we see that $\text{Trace}_{L/K}$ is identically zero.

\medskip\noindent
Assume on the other hand that $L/K$ is separable.
By induction on the degree we will show that
$\text{Trace}_{L/K}$ is not identically zero.
Thus by Lemma \ref{lemma-trace-and-norm-tower} we may assume that
$L/K$ is generated by a single element $\alpha$ (use that if
the trace is nonzero then it is surjective).
We have to show that $\text{Trace}_{L/K}(\alpha^e)$
is nonzero for some $e \geq 0$.
Let $P = x^d + a_1 x^{d - 1} + \ldots + a_d$ be the minimal
polynomial of $\alpha$ over $K$.
Then $P$ is also the characteristic polynomial of the linear
maps $\alpha : L \to L$, see
Lemma \ref{lemma-characteristic-vs-minimal-polynomial}.
Since $L/k$ is separable we see from Lemma \ref{lemma-recognize-separable}
that $P$ has $d$ pairwise distinct roots $\alpha_1, \ldots, \alpha_d$
in an algebraic closure $\overline{K}$ of $K$. Thus these are the
eigenvalues of $\alpha : L \to L$.
By linear algebra, the trace of $\alpha^e$ is
equal to $\alpha_1^e + \ldots + \alpha_d^e$.
Thus we conclude by Lemma \ref{lemma-sums-of-powers}.
\end{proof}

\noindent
Let $K$ be a field and let $Q : V \times V \to K$ be a bilinear form
on a finite dimensional vector space over $K$. Say $\dim_K(V) = n$.
Then $Q$ defines a linear map $Q : V \to V^*$, $v \mapsto Q(v, -)$
where $V^* = \Hom_K(V, K)$ is the dual vector space. Hence a linear map
$$
\det(Q) : \wedge^n(V) \longrightarrow \wedge^n(V)^*
$$
If we pick a basis element $\omega \in \wedge^n(V)$, then we can
write $\det(Q)(\omega) = \lambda \omega^*$, where $\omega^*$
is the dual basis element in $\wedge^n(V)^*$. If we change our
choice of $\omega$ into $c \omega$ for some $c \in K^*$, then
$\omega^*$ changes into $c^{-1} \omega^*$ and therefore
$\lambda$ changes into $c^2 \lambda$. Thus the class of
$\lambda$ in $K/(K^*)^2$ is well defined and is called the
{\it discriminant of $Q$}. Unwinding the definitions we see that
$$
\lambda = \det(Q(v_i, v_j)_{1 \leq i, j \leq n})
$$
if $\{v_1, \ldots, v_n\}$ is a basis for $V$ over $K$. Observe that
the discriminant is nonzero if and only if $Q$ is nondegenerate.

\begin{definition}
\label{definition-discriminant}
Let $L/K$ be a finite extension of fields. The
{\it discriminant of $L/K$} is the discriminant of
the trace pairing $Q_{L/K}$.
\end{definition}

\noindent
By the discussion above and Lemma \ref{lemma-separable-trace-pairing}
we see that the discriminant is nonzero if and only
if $L/K$ is separable. For $a \in K$ we often say
``the discriminant is $a$'' when it would be more correct
to say the discriminant is the class of $a$ in $K/(K^*)^2$.

\begin{exercise}
\label{exercise-quadratic-discriminant}
Let $L/K$ be an extension of degree $2$. Show that exactly
one of the following happens
\begin{enumerate}
\item the discriminant is $0$, the characteristic of $K$ is $2$,
and $L/K$ is purely inseparable obtained by taking a square root
of an element of $K$,
\item the discriminant is $1$, the characteristic of $K$ is $2$, and
$L/K$ is separable of degree $2$,
\item the discriminant is not a square, the characteristic of $K$
is not $2$, and $L$ is obtained from $K$ by taking the square root
of the discriminant.
\end{enumerate}
\end{exercise}











